%%%%%%%%%%%%%%%%%%%%%%%%%%%%%%%%%%%%%%%%%%%%%%%%
% Ausleihliste Werkzeug (Formular zum Ausdrucken)
% CC-BY-SA 3.0, FAU FabLab
% Früher ausleihliste/AusleihlisteWerkzeug.xls (LibreOffice), jetzt LaTeX.
%%%%%%%%%%%%%%%%%%%%%%%%%%%%%%%%%%%%%%%%%%%%%%%%
\documentclass[a4paper,landscape,10pt]{scrartcl}
\usepackage{fablab-aushang}

% ---------------------------------------------------------------------
%  Spalten der Liste: Breite und Überschrift
% ---------------------------------------------------------------------
\newcommand{\hA}{10}      % Zeile zum Ausfüllen (mm)
\newcommand{\hK}{5.5}     % Kommentarzeile für das Pfand (mm)
\newcommand{\Eintraege}{7}

\begin{document}

% ------------------------------ Kopf ---------------------------------
\Kopf{Ausleihliste Werkzeug}{}
\begin{minipage}[t]{0.70\linewidth}
  \textbf{Werkzeuge oder Geräte werden nur verliehen, wenn}
  \begin{itemize}\setlength{\itemsep}{0pt}\setlength{\parskip}{0pt}
    \item es einen guten Grund gibt, dass du es außerhalb des FabLabs brauchst,
    \item und es währenddessen voraussichtlich nicht im FabLab gebraucht wird,
    \item und ein Betreuer die Erlaubnis erteilt,
    \item und ein Betreuer für dich bürgt oder du ein Pfand hinterlässt
          (z.\,B. Führerschein, FAU-Karte, Geld -- \textbf{kein Personalausweis oder Pass!}).
  \end{itemize}
\end{minipage}\hfill
\begin{minipage}[t]{0.27\linewidth}
  \fcolorbox{fablab}{kopf}{\parbox{\dimexpr\linewidth-2\fboxsep-2\fboxrule\relax}{%
    \textbf{Bei der Rückgabe:}\\
    Werkzeug auf Beschädigung und Vollständigkeit prüfen.}}
\end{minipage}

\vspace{3mm}
% ------------------------------ Liste --------------------------------
\begin{tikzpicture}[x=1mm, y=1mm, font=\small]
  % Spalten: linke Kante und Breite in mm (zusammen 273 mm)
  \def\spalten{0/20/Datum,
               20/22/Rückgabe\\spätestens,
               42/40/Werkzeug,
               82/38/{Vorname,\\Nachname},
               120/34/{Unterschrift\\des/der Entleihenden},
               154/44/{Kontaktdaten\\(E-Mail, Handynummer)},
               198/32/{Name und\\Unterschrift\\des Betreuers},
               230/20/Datum,
               250/23/Unterschrift\\des Betreuers\\(leserlich)}
  \def\breite{273}
  \def\ausleihe{230}   % Ende der Spalten für die Ausleihe
  % Gruppenüberschriften
  \fill[kopf] (0,0) rectangle (\breite,-6);
  \node[font=\bfseries] at (0.5*\ausleihe, -3) {Ausleihe};
  \node[font=\bfseries] at (0.5*\ausleihe+0.5*\breite, -3) {Rückgabe};
  % Spaltenüberschriften
  \foreach \x/\w/\t in \spalten {
    \node[anchor=north west, align=left, font=\footnotesize\bfseries, text width=\w mm-2mm]
      at (\x+1, -7) {\t};
  }
  \pgfmathsetmacro{\yk}{-19}   % Unterkante des Tabellenkopfs
  \draw[line width=0.8pt] (0,-6) -- (\breite,-6);
  \draw[line width=0.8pt] (0,\yk) -- (\breite,\yk);
  % Einträge
  \foreach \i in {1,...,\Eintraege} {
    \pgfmathsetmacro{\yo}{\yk-(\i-1)*(\hA+\hK)}
    \pgfmathsetmacro{\ym}{\yo-\hA}
    \pgfmathsetmacro{\yu}{\yo-\hA-\hK}
    % Zeile zum Ausfüllen
    \foreach \x/\w/\t in \spalten {
      \ifdim\x pt<\ausleihe pt\relax
        \draw[black!60, line width=0.3pt] (\x,\yo) -- (\x,\ym);
      \fi
    }
    \draw[black!60, line width=0.3pt] (0,\ym) -- (\ausleihe,\ym);
    % Kommentarzeile über alle Ausleihe-Spalten
    \node[anchor=west, font=\footnotesize] at (1,\ym-0.5*\hK) {Kommentar (Pfand):};
    % Rückgabe über beide Zeilen
    \draw[black!60, line width=0.3pt] (250,\yo) -- (250,\yu);
    \draw[line width=0.8pt] (0,\yu) -- (\breite,\yu);
  }
  % Rahmen und Trennlinie Ausleihe/Rückgabe
  \pgfmathsetmacro{\yende}{\yk-\Eintraege*(\hA+\hK)}
  \draw[line width=1.2pt] (0,0) rectangle (\breite,\yende);
  \draw[line width=1.2pt] (\ausleihe,0) -- (\ausleihe,\yende);
  \foreach \x/\w/\t in \spalten {
    \ifdim\x pt>0pt\relax
      \draw[black!60, line width=0.3pt] (\x,-6) -- (\x,\yk);
    \fi
  }
\end{tikzpicture}

\end{document}
